%% ============================================================
%% 分片 04：原文第 3 节「Topological Field Theory From ${\EUN}=4$ Super Yang-Mills Theory」（原文 glshort.tex 起于第 1216 行）
%% 由 tools/harvmac2tex.py 从 arXiv 源码（hep-th/0604151v3）自动生成。
%% 只译英文散文；公式、编号、\label、\cite 一律不改。
%% ============================================================

\section{从 ${\EUN}=4$ 超对称杨--米尔斯理论得到的拓扑场论} \label{sec:tqftfour}

\begin{keytakeaways}
\begin{itemize}[leftmargin=1.4em,itemsep=2pt,topsep=2pt]
\item \textbf{物理图像}：把 ${\EUN}=4$ 超对称杨--米尔斯理论的自旋群 $Spin(4)$ 换成它在 $Spin(4)\times Spin(6)$ 中的对角子群——这就是扭化；其效果是让某个超荷变成标量荷 $Q$，从而在任意四维流形上得到一族以 $\Bbb{CP}^1$ 为参数的拓扑场论，$S$-对偶作用在这个参数上。
\item \textbf{数学工具}：$Spin(4)$ 与 $R$-对称群 $Spin(6)=SU(4)_R$ 的表示论、${\EUN}=4$ 超对称变换律、十维旋量在四维语言下的分解、标量超荷确定的 $Q$-上同调，以及后文的消灭定理与单圈局域化。
\item \textbf{关键结论}：扭化后的拓扑超荷取 $Q=uQ_\ell+vQ_r$，参数 $t=v/u$ 只在 $t\mapsto-t$ 下等价（这一等价来自经典作用量的 $U(1)$ 对称性）；因此本质的区分是按 $t=\pm i$（平坦 $G_{\Bbb{C}}$ 联络）与实 $t$（椭圆方程、计数解）来划分的。
\item \textbf{边界与展望}：扭化在 Lorentz 号差下并不好用，故全文以欧氏号差为准，只在 $M=\Bbb{R}\times W$、$\Bbb{R}^2\times C$ 这类有「时间方向」的分解下临时回到 Lorentz 号差。本节随后将写出超对称拉格朗日量，并引入典范参数 $\Psi$。
\end{itemize}
\end{keytakeaways}

我们的下一个目标是描述如何通过扭化 ${\EUN}=4$ 超对称杨--米尔斯理论，在任意四维流形\footnote{我们通常假定 $M$ 可定向，因为此时构造更有趣。然而，如果只限于式~\eqref{eq:oopy} 的 $\theta$ 角为零、且第~\ref{sec:canonpar}~节的典范参数 $\Psi$ 等于 $0$ 或 $\infty$ 的情形，那么即使 $M$ 不可定向，该构造也有意义。} $M$ 上构造四维拓扑量子场论（TQFT）。在构造 TQFT 时，使用欧氏号差最为自然：拓扑场论与欧氏号差之间的联系本就最为自然，而且无论如何，我们所用的扭化在 Lorentz 号差下都不太好用。当我们取 $M=\Bbb{R}\times W$ 或 $M=\Bbb{R}^2\times C$（$W$ 为三维流形、$C$ 为二维流形）时，让 $\Bbb{R}$ 或 $\Bbb{R}^2$ 取 Lorentz 号差，就能有用地回到 Lorentz 号差。

\begin{yizhu}[物理图像]
按：脚注里「$\theta=0$ 且典范参数 $\Psi=0$ 或 $\infty$ 时构造对不可定向 $M$ 也有意义」这一句，值得记下它的来源。四维 TQFT 的构造通常依赖定向（因为要与体积形式配对）；而 $\theta$ 项 $\theta\int_M\Tr F\wedge F$ 中的被积量就是瞬子数密度，它在定向反转下变号——相应地，定向反转在扭化参数上的作用为 $\Psi\mapsto-\Psi$。因此若 $\theta=0$ 且 $\Psi$ 取端点值 $0$ 或 $\infty$（它们在 $\Psi\mapsto-\Psi$ 下不动），定向反转就不改变理论，构造便可延拓到不可定向的四维流形。关于扭化与 Hitchin 方程的一般图景，可参见 Witten 2008 综述。
\end{yizhu}

转到欧氏号差会带来一个重要变化：$Spin(4)$ 的两个自旋表示 $(\btwo,\bone)$ 与 $(\bone,\btwo)$ 是伪实的，而在 $Spin(1,3)$ 下这两个表示互为复共轭。相应地，区分这两个自旋表示的算子 $\Gamma_0\Gamma_1\Gamma_2\Gamma_3$ 现在平方为 $+1$；为避免混淆，我们在使用欧氏号差时称它为 $\hat\Gamma_E$，而把 $\hat\Gamma$ 留给 Lorentz 号差的情形。类似地，两个自旋表示 ${\EUS}^\pm$ 在欧氏号差下互为复共轭，区分它们的算子 $\Gamma_0\Gamma_1\dots\Gamma_9$——我们现在称它为 $\bar\Gamma_E$——在 ${\EUS}^\pm$ 上以 $\mp i$ 为本征值。

\begin{yizhu}[易错点]
按：这一段同时出现两个「手征算子」，极易混淆。$\hat\Gamma_E=\Gamma_0\Gamma_1\Gamma_2\Gamma_3$ 是四维的，它区分 $Spin(4)$ 的两个 Weyl 表示 $(\btwo,\bone)$ 与 $(\bone,\btwo)$；而 $\bar\Gamma_E=\Gamma_0\cdots\Gamma_9$ 是十维的，它区分十维旋量的手征性，在四维语言里就是区分 ${\EUS}^\pm$。式~\eqref{eq:nohinks} 里的 $-i$ 只是一个约定，取决于我们选定的 Wick 转动方向 $ix^0_L=x^0_E$；但它是全文符号簿记的起点：后文式~\eqref{eq:banother} 的 $\Gamma_{89}\epsilon=-i\epsilon$、以及扭化超荷的 $\EUK=-1$ 标定，都由此决定。若把转动方向反过来，全篇的 $i$ 会一致地变号，而物理结论不变。
\end{yizhu}

为确定符号，约定下标 $L$ 或 $E$ 分别指 Lorentz 号差或欧氏号差；我们通过 $ix^0_L=x^0_E$ 做从 Lorentz 号差到欧氏号差的 Wick 转动，于是 $\partial/\partial x^0_L=i\partial/\partial x^0_E$，$\Gamma_{0\,L}=i\Gamma_{0\,E}$。因此 $\bar\Gamma_L=i\bar\Gamma_E$；又因 $\bar\Gamma=\bar\Gamma_L$ 在 ${\EUS}^\pm$ 上以乘以 $\pm 1$ 的方式作用，$\bar\Gamma_E$ 便以乘以 $\mp i$ 的方式作用。特别地，若 $\epsilon$ 是超对称生成元，则
% 式 (3.1)
\setcounter{equation}{0}
\begin{equation}\label{eq:nohinks}
\bar\Gamma_E\epsilon=-i\epsilon.
\end{equation}

\begin{yizhu}[含糊点拨]
按：原文说「扭化在 Lorentz 号差下不好用」，用本段的表示论语言就能看清原因。扭化要求把四维自旋群的自旋表示与 $R$-对称群 $Spin(6)$ 的表示放在同一个实形式上配对；在 Lorentz 号差下 $Spin(1,3)$ 的两个 Weyl 表示互为复共轭，十维旋量的手征性与四维的左右手性无法像欧氏号差下那样统一处理——本段所说的「伪实」正是欧氏号差才具有的性质。因此全文以欧氏号差为准，只在 $M=\Bbb{R}\times W$、$\Bbb{R}^2\times C$ 这类分解下临时回到 Lorentz 号差。
\end{yizhu}
